\documentclass[12pt]{article}
\usepackage{amsmath, amssymb}
\usepackage{geometry}
\geometry{margin=1in}

\begin{document}

\section*{CSE400 – Fundamentals of Probability in Computing}

\subsection*{Lecture 4: Joint Probability and Conditional Probability}

\textbf{Instructor:} Dhaval Patel, PhD\\
\textbf{Date:} January 15, 2026

\hrule
\vspace{1em}

\section*{1. Motivation and Engineering Applications}

\subsection*{1.1 Purpose of Learning Probability}

The lecture motivates probability theory through engineering systems where outcomes are uncertain and must be represented using likelihoods.

\subsection*{1.2 Engineering Application 1: Speech Recognition System}

\subsubsection*{Vocabulary and Templates}

\begin{itemize}
  \item Vocabulary shown:
  \begin{itemize}
    \item Hello
    \item Yes
    \item No
    \item Bye
  \end{itemize}
  \item Each word corresponds to a \textbf{template waveform}.
\end{itemize}

\subsubsection*{Speaker Variability}

\begin{itemize}
  \item The same word spoken by different speakers produces different waveforms:
  \begin{itemize}
    \item Speaker 1 (male)
    \item Speaker 2 (male)
    \item Speaker 3 (female)
    \item Speaker 4 (child)
  \end{itemize}
\end{itemize}

\subsubsection*{Signal Representation}

\begin{itemize}
  \item Time-domain signal: $( d(t) )$
  \item Transformation shown to frequency domain:
  \[
  d(t) \rightarrow X(\omega)
  \]
\end{itemize}

\subsubsection*{Noise and Interference}

\begin{itemize}
  \item Observed signal includes:
  \begin{itemize}
    \item Noise
    \item Interference
  \end{itemize}
\end{itemize}

\subsubsection*{Probability Emphasis}

\begin{itemize}
  \item A probability comparison is written:
  \[
  P_2 > P_1
  \]
  \item Recognition is based on which outcome is \textbf{more probable}.
\end{itemize}

\subsection*{1.3 Engineering Application 2: Radar System}

\subsubsection*{System Description}

\begin{itemize}
  \item Transmitter (Tx) sends a signal.
  \item Receiver (Rx) observes the signal.
\end{itemize}

\subsubsection*{Hypotheses Shown}

\[
H_0: y_i = w_i
\]
\[
H_1: y_i = s_i + w_i
\]

\begin{itemize}
  \item $( s_i )$: signal
  \item $( w_i )$: noise
\end{itemize}

\subsubsection*{Detection Outcomes}

\begin{itemize}
  \item False Alarm
  \item Miss Detection
\end{itemize}

\subsubsection*{Probability Relationship}

\[
P_m + P_d = 1
\]

\begin{itemize}
  \item Noise is labeled as \textbf{Gaussian} on the slide.
\end{itemize}

\subsection*{1.4 Engineering Application 3: Communication Network}

\subsubsection*{Network Description}

\begin{itemize}
  \item Multiple nodes from \textbf{source} to \textbf{destination}.
  \item Packets traverse different paths.
\end{itemize}

\subsubsection*{Random Quantities}

\begin{itemize}
  \item Arrival of packets
  \item Delay (latency)
\end{itemize}

\subsubsection*{Quality of Service (QoS)}

\begin{itemize}
  \item Delay shown in:
  \begin{itemize}
    \item milliseconds (ms)
    \item microseconds ($\mu$s)
  \end{itemize}
\end{itemize}

\subsubsection*{Wireless Technologies Listed}

\begin{itemize}
  \item WiFi:
  \begin{itemize}
    \item 2.4 GHz
    \item 5 GHz
    \item 6 GHz
  \end{itemize}
\end{itemize}

\section*{2. Lecture Outline (As Shown)}

\subsection*{2.1 Topics Listed}

\begin{enumerate}
  \item \textbf{Introduction to Probability Theory}
  \begin{itemize}
    \item Experiments, Sample Space, and Events
    \item Axioms of Probability
    \item Corollaries and Propositions from Probability Axioms
    \item How to Assign Probability
  \end{itemize}

  \item \textbf{Joint Probability}
  \begin{itemize}
    \item Motivation, Notation, and Concepts
    \item Example 1: Card Deck Example
    \item Example 2: Costume Party Example
  \end{itemize}

  \item \textbf{Conditional Probability}
  \begin{itemize}
    \item Motivation, Notation, and Concepts
    \item Example 3: Cards Without Replacement
    \item Example 4: Game of Poker
    \item Example 5: The Missing Key
  \end{itemize}
\end{enumerate}

\begin{quote}
Only \textbf{Introduction to Probability Theory} is developed in the provided slides.\\
Joint Probability and Conditional Probability are \textbf{not expanded further} and are therefore \textbf{not included beyond this outline}.
\end{quote}

\section*{3. Introduction to Probability Theory}

\subsection*{3.1 Experiments, Sample Space, and Events}

\section*{4. Definitions}

\subsection*{4.1 Experiment}

\textbf{Definition:}\\
An \textbf{Experiment (E)} is a procedure we perform that produces some result.

\textbf{Example:}
\begin{itemize}
  \item Tossing a coin five times
  \item Denoted as $( E_5 )$
\end{itemize}

\subsection*{4.2 Outcome}

\textbf{Definition:}\\
An \textbf{Outcome $(\xi)$} is a possible result of an experiment.

\textbf{Example:}
\[
\xi_1 = \text{HHTHT}
\]

\subsection*{4.3 Event}

\textbf{Definition:}\\
An \textbf{Event (any letter)} is a certain set of outcomes of an experiment.

\textbf{Example:}
\[
C = \{\text{all outcomes consisting of an even number of heads}\}
\]

\subsection*{4.4 Sample Space}

\textbf{Definition:}\\
The \textbf{Sample Space (S)} is the collection or set of \textbf{all possible distinct outcomes} of an experiment.

\textbf{Properties Listed}
\begin{itemize}
  \item Mutually Exclusive
  \item Collectively Exhaustive
\end{itemize}

\textbf{Example: Flip a Coin}
\begin{itemize}
  \item Mutually Exclusive: Heads or tails, not both.
  \item Collectively Exhaustive: No outcome other than heads or tails.
\end{itemize}

\subsection*{4.5 Sample Space (Continued)}

\begin{itemize}
  \item $( S )$ is the \textbf{universal set} of outcomes of an experiment.
  \item Sample space types:
  \begin{itemize}
    \item Discrete
    \item Countably infinite
    \item Continuous
  \end{itemize}
\end{itemize}

\subsection*{4.6 Examples of Experiments (Listed)}

\begin{enumerate}
  \item Flipping a fair coin once
  \item Cubical die with numbered faces is rolled
  \item Rolling two dices
  \item Flip a coin until tails occurs
  \item Random number generator with interval $([0,1))$
\end{enumerate}

\section*{5. Axioms of Probability}

\subsection*{5.1 Probability}

\textbf{Definition:}\\
Probability is:
\begin{itemize}
  \item A measure of the likelihood of various events\\
  \textbf{OR}
  \item A function of an event that produces a numerical quantity that measures the likelihood of that event.
\end{itemize}

It is stated that:
\begin{itemize}
  \item There are many ways to define such a function.
\end{itemize}

\subsection*{5.2 Axioms}

\textbf{Statement:}\\
The axioms are self-evident and require no proof.

\subsection*{5.3 Axiom 1}

For any event $( A )$:
\[
0 \leq \Pr(A) \leq 1
\]

\subsection*{5.4 Axiom 2}

If $( S )$ is the sample space:
\[
\Pr(S) = 1
\]

\subsection*{5.5 Axiom 3}

\subsubsection*{(a) Finite Mutually Exclusive Events}

If $( A \cap B = \varnothing )$:
\[
\Pr(A \cup B) = \Pr(A) + \Pr(B)
\]

\subsubsection*{(b) Infinite Mutually Exclusive Events}

For mutually exclusive sets $( A_i )$, $( i = 1,2,3,\ldots )$:
\[
\begin{array}{c}
\Pr\left(\bigcup_{i=1}^{\infty} A_i\right) \\
\hline
\sum_{i=1}^{\infty} \Pr(A_i)
\end{array}
\]

\section*{6. Corollary from Probability Axioms}

\subsection*{6.1 Corollary 2.1}

\textbf{Definition:}\\
Define $( A_i )$ as the null event for all values of $( i )$ greater than $( n )$.

\textbf{Statement:}\\
For \textbf{M finite} mutually exclusive sets $( A_i )$:
\[
\begin{array}{c}
\Pr\left(\bigcup_{i=1}^{M} A_i\right) \\
\hline
\sum_{i=1}^{M} \Pr(A_i)
\end{array}
\]

\subsection*{6.2 Relation to Axiom 3}

\begin{itemize}
  \item Axiom 3 is equivalent to Corollary 2.1 when the sample space is finite.
  \item The generality of Axiom 3 is necessary when the sample space has infinitely many points.
\end{itemize}

\section*{7. Propositions from Probability Axioms}

\subsection*{7.1 Proposition 2.1}

\[
\Pr(A^c) = 1 - \Pr(A)
\]

\subsection*{7.2 Proposition 2.2}

If $( A \subset B )$:
\[
\Pr(A) \leq \Pr(B)
\]

\section*{8. End of Lecture Content}

\end{document}
